% ============================================================================
\section{Limitations and Conclusion}
\label{sec:conclusion}
% ============================================================================

The study evaluates one student scale, four tasks, and a serial one-teacher-slot
system. GPAS is a local rule based on the current within-task noise, while the
analysis treats equal-size micro-batches as independent and holds the AdamW
scaling fixed within a step. Nonrepeating prompt streams support the sampling
assumption but cannot remove correlations induced by a shared student. The
additive timing model is specific to serial execution; concurrent teacher slots
require a different Cost-GPAS objective. Inverse-initial-loss weighting is also
unstable when an initial task KL is close to zero. When the measured
variance ratio $H$ remains near one, the predicted variance under GPAS is close
to Uniform, and Uniform is simpler to use.

GPAS fixes the objective weights while adapting the precision of each task's
gradient estimate. Scaling by $w_i/m_i$ preserves the expected gradient. GPAS
uses optimizer-scaled noise to set the micro-batch counts, and Cost-GPAS also
uses measured task time. Both methods use the gradients and timings produced by
standard MOPD training and work with conventional AdamW. The taskwise
second-moment correction is optional. \missing{Replace this sentence with the
main four-teacher empirical conclusion once the results are available.}
